\chapter{Lifecycle, Tasks, and Cancellation}
\label{ch:lifecycle-tasks}
\chapterquote{The object that installs a callback should also know how that callback dies.}{The cleanup law}

\begin{chaptermap}
Core question & How does a Swing application avoid stale callbacks, duplicate listeners, unreleased caches, and obsolete background results?\\
Primary APIs & \api{Disposable}, \api{Subscription}, \api{SubscriptionScope}, \api{Detachable}, \api{DetachableScope}, \api{TaskService}, \api{TaskHandle}, \api{TaskCallbacks}, \api{CancellationToken}, \api{GenerationCounter}, \api{AsyncFieldValidation}, and \api{SwingTaskServices}.\\
Plain Swing baseline & Swing exposes listeners, timers, workers, models, and callbacks; it does not automatically give them one lifetime owner.\\
Swing connection & Binding scopes, behavior scopes, task callbacks, busy overlays, dialog disposal, view activation, and stale-result suppression on the EDT.\\
Companion examples & \lifecycleExample, \detachableScopeExample, \asyncTaskExample, \generationCounterExample, \swingTaskExample, and \bookAppExample.\\
\end{chaptermap}

\section{Lifecycle is architecture}

Every nontrivial Swing screen installs things. It installs action listeners,
document listeners, property-change listeners, list-selection listeners,
observable-value subscriptions, table controllers, timers, keyboard actions,
validation decorators, task callbacks, overlays, and service observers. A
screen may be opened and closed hundreds of times during a long application
session. If the relationships installed by the screen survive the screen that
needed them, the application leaks memory, behavior, or both.

\termdef{Lifecycle}{in this chapter} means the interval during which an object
is allowed to receive callbacks and own resources. A modal dialog's lifecycle
may be one call to show the dialog. A workspace panel's lifecycle may be one
activation inside a tab. A presenter may live longer than its current Swing
view. A cached loadable model may survive deactivation but detach its cached
data. Lifecycle is not a synonym for Java object reachability. It is the
program's declared ownership boundary.

\termdef{Callback}{in this discussion} is any function, listener, observer, or
action that another object may call later. Swing is callback-rich: listeners
are the ordinary way that components, documents, models, timers, and workers
communicate. Callback-rich code is not automatically leak-prone. It becomes
leak-prone when the callback source lives longer than the object captured by
the callback. A document listener that captures a panel keeps the panel alive
as long as the document keeps the listener. A task callback that captures a
dialog may update the dialog after the user has closed it. A timer that
captures a component can keep repainting a dead view.

The rule is simple enough to memorize:

\begin{principlebox}{the cleanup law}
The object that installs a callback should either close it directly or register
it in a scope that closes with the same lifecycle.
\end{principlebox}

The rule is deliberately stronger than ``remember to remove listeners.'' It
asks for a visible owner. A line of code installs a relationship; a nearby line
registers the cleanup handle; a lifecycle boundary closes the scope. This turns
cleanup from a hope into a reviewable fact. A reviewer can ask ``what closes
this?'' and expect the answer to be a concrete object, not a story about the
garbage collector.

\begin{figure}[htbp]
\centering
\begin{minipage}[t]{0.48\linewidth}
\includegraphics[width=\linewidth]{\busyTaskFigure}
\end{minipage}\hfill
\begin{minipage}[t]{0.48\linewidth}
\includegraphics[width=\linewidth]{\backgroundStatusFigure}
\end{minipage}
\caption{Lifecycle and task ownership become visible through busy overlays, progress, cancellation controls, and status text. Each visual relationship needs a cleanup owner.}
\label{fig:lifecycle-task-screens}
\end{figure}

Figure~\ref{fig:lifecycle-task-screens} shows visual states, but the states are
not only visual. A busy overlay depends on a busy value. A progress bar depends
on task progress or task state. A Cancel button depends on a task handle. A
status label depends on callbacks. Each dependency is a relationship that must
end. If the view closes while the task is still running, the task may continue
or be cancelled according to policy, but it must not retain the closed view as
a callback target.

\begin{pitfallbox}{closed is not collected}
Calling \api{dispose} on a window removes the native peer. It does not
magically remove every listener, subscription, task callback, timer, or service
reference that points at the window's object graph.
\end{pitfallbox}

Swing leaks often announce themselves first as behavior, not heap usage. A
screen is opened, closed, and opened again; validation messages appear twice.
A search field changes; two requests are sent for each edit. A closed dialog
briefly updates a status label when an old task completes. A timer continues to
repaint a panel that no longer has a native peer. Later the heap confirms the
same story: old panels remain reachable through listener lists, task handles,
or service callbacks. The bug was already visible in the ownership graph.

It is useful to list the common lifecycle boundaries before code is written.
Application lifetime covers things such as resource bundles, service
registries, logging, and top-level command registries. Window lifetime covers
native shell resources and window-specific global actions. Screen activation
covers bindings between a presenter and a view. Dialog activation covers
temporary form bindings, keyboard policies, task services, and validation
decorations. Selection lifetime covers details loaded for one selected master
row. Request lifetime covers one background lookup or validation request.
Different boundaries may overlap; they should not be confused.

The dangerous case is a long-lived object that stores a callback capturing a
short-lived object. A repository event bus stores a listener that captures a
customer panel. A global command registry stores an action that captures a
dialog. A worker service stores a callback that captures a table model owned by
a closed tab. A property-change listener on a shared model captures a
presenter. In each case the source is legitimate and may live for the whole
application. The captured object does not. The cleanup law asks us to register
the removal when the relationship is installed.

Weak listeners are not a general answer. They can help when an external API
forces a long-lived listener source and the listener is purely observational.
They are not a replacement for deterministic close handles. A weak listener
that disappears too early loses behavior. A weak listener that remains because
some other object still references it still leaks. A weak listener also does
not solve stale task results. Prefer deterministic removal when the lifecycle
boundary is known, which is most of the time in ordinary UI code.

A useful review device is the ownership ledger. For each screen, make a small
table with three columns: installed relationship, owner, and close boundary.
The rows are ordinary: document listener, selection listener, table model
adapter, command adapter, validation binding, task service, task handle, timer,
overlay binding, generated behavior scope, and detachable cache. The owner
column names the object or scope that closes it. The close boundary names the
event that ends ownership: dialog close, tab deactivation, workspace close,
application shutdown, or selection change. The ledger is not meant to become
heavy documentation. It is a way to discover that three relationships have no
owner before they become defects.

Lifecycle smells are easy to recognize after using the ledger for a while. A
screen has a field named \api{closed} that many callbacks check, but the
callbacks are never removed. A task callback checks whether a component is
displayable before updating it, but the task still retains the component until
completion. A presenter has a \api{dispose} method that closes some listeners
while others are removed in anonymous window listeners. A global event bus has
registrations in constructors but no corresponding unregister calls. A dialog
creates a new task service for each Apply click and never closes any of them.
Each smell says the same thing: a relationship was installed without an owner
whose lifetime matches the relationship.

The ownership boundary should also match the user's mental model. When a user
closes a dialog, work that exists only for that dialog should stop or become
irrelevant. When a user switches master selection, the previous detail load
should stop or become irrelevant. When a user hides a tab temporarily, cached
data may be released but the tab's durable settings may remain. When the
application shuts down, user prompts and cleanup may follow a different policy
from ordinary close. The code should make these differences visible instead of
using one \api{disposeEverything} method for every boundary.

\begin{coruscobox}{What Corusco changes}
Corusco gives callback ownership small, reusable names. A listener attachment
can return a \api{Subscription}; a set of cleanup handles can live in a
\api{SubscriptionScope}; reusable cached state can implement \api{Detachable};
background work can return a \api{TaskHandle}. The names are plain because they
must be used everywhere.
\end{coruscobox}

\section{Deterministic cleanup with scopes}

\termdef{Disposable}{in Corusco} is the small close contract. It extends
\api{AutoCloseable} with an unchecked close signature suitable for UI cleanup
handles. A disposable may remove a listener, cancel a task, restore a border,
clear a keyboard binding, release a cache, or close another scope. The caller
does not need to know the original removal API; the disposable handle knows how
to undo the relationship it represents.

\termdef{Subscription}{in Corusco} is a disposable registration for a listener,
observer, or similar attachment. Observable values and lists return
subscriptions when code subscribes to change events. Swing bindings and
behavior installations often return disposable handles with the same ownership
shape. The exact name matters less than the rule: if a call installs a callback,
the call should return or register something that can remove it.

\begin{lstlisting}[caption={A subscription is a cleanup handle.}]
Subscription subscription = value.subscribe(event -> updateStatus(event));

// Later, when this owner is no longer active:
subscription.close();
\end{lstlisting}

\api{SubscriptionScope} owns a group of disposables and closes them
deterministically. It is intended for a single ownership context such as a
presenter activation, dialog activation, behavior, or test fixture. Children
close in reverse registration order. Closing the scope is idempotent. If one
child fails, the scope still attempts to close the remaining children and then
throws a \api{SubscriptionScopeException} with the individual failures
suppressed. If a child is added after the scope has closed, the child is closed
immediately and is not retained.

\begin{lstlisting}[caption={A screen activation owns one cleanup scope.}]
final class CustomerPanelController implements AutoCloseable {
    private final SubscriptionScope scope = new SubscriptionScope();

    CustomerPanelController(CustomerPresenter presenter, CustomerView view) {
        scope.add(presenter.statusText().subscribe(event ->
                view.statusLabel().setText(event.newValue())));

        scope.add(view.refreshButtonBinding());
        scope.onClose(() -> System.out.println("customer panel closed"));
    }

    @Override
    public void close() {
        scope.close();
    }
}
\end{lstlisting}

The example is small, but it changes review. A reviewer can ask where the
scope closes. If the answer is visible, listener cleanup is not folklore. The
constructor installs relationships and immediately registers their cleanup
handles. The controller's close method closes the scope. The object that owns
the controller must then close the controller at the screen boundary. Each
layer has one job.

Reverse close order is not cosmetic. Relationships are often built in layers.
A source value is observed by a derived value. The derived value is observed by
a validation binding. The binding decorates a component. The component is
inside a dialog whose keyboard binding is registered last. Closing in reverse
order tears down the outermost relationship first, then the relationships it
depended on. This mirrors stack discipline. It should not be abused for
ordinary business logic, but it reduces surprises when resources are nested.

Fail-closed late registration is just as important. During teardown, callbacks
may still arrive. A callback may create a helper, subscribe to a value, and try
to register it with a scope that has already closed. If the scope silently
accepts the child, the application has created a new leak while closing the old
one. \api{SubscriptionScope.add} closes late children immediately. This makes
the safe outcome the default. The late child does not run because the owner is
already dead.

\begin{lstlisting}[caption={Late registration fails closed.}]
scope.close();

Subscription late = scope.add(value.subscribe(event -> updateClosedView()));
assertThat(scope.isClosed()).isTrue();
\end{lstlisting}

The late subscription is returned for fluency, but it is not retained by the
scope. Its cleanup has already run. A test can assert the source listener count
if the source exposes one, or it can mutate the source and assert that the
closed view does not update. This small behavior is a powerful guard against
race conditions around deactivation.

Failure during cleanup should be visible but not prevent other cleanup. If one
child close throws and the scope stopped immediately, later listeners would
remain installed because an earlier cleanup was buggy. Corusco scopes instead
continue closing the rest and collect failures. This follows the same principle
as structured cleanup in systems code: attempt to release everything, then
report what failed. UI code should not throw during close casually, but if it
does, other relationships still deserve cleanup.

Scopes are not synchronized. A \api{SubscriptionScope} is usually owned and
closed on one UI or presenter thread. If a project shares a scope across
threads, it must provide coordination. In Swing code, the simpler rule is to
create, add, and close UI scopes on the EDT. Background workers may hold task
handles, but they should not mutate Swing-owned scopes directly. Use a callback
executor or \api{SwingEdt.executor} to return ownership-sensitive work to the
EDT.

Scopes compose with Java's \api{try}-with-resources when the lexical lifetime
matches the logical lifetime. Tests and small command blocks often benefit:
create a scope, install relationships, assert behavior, and let the scope close
at the end of the block. Long-lived screens usually close scopes from explicit
deactivation or disposal methods instead. The important point is not the Java
syntax; it is that the lifetime is visible.

\begin{lstlisting}[caption={A small test can use lexical scope ownership.}]
try (SubscriptionScope scope = new SubscriptionScope()) {
    scope.add(model.value().subscribe(event -> seen.add(event.newValue())));
    model.setValue("updated");
}

model.setValue("late");
assertThat(seen).containsExactly("updated");
\end{lstlisting}

Corusco Swing bindings and behavior installations build on the same ownership
idea. A direct value binding, validation tooltip, red-border behavior, composed
tooltip, command adapter, or busy overlay binding should return a handle or be
registered in a scope. A generated behavior plan should install into a
behavior scope that closes with the view or dialog activation. The Swing
object may look visual, but its listener relationships are lifecycle facts.

\begin{coruscobox}{Scopes as binding owners}
Corusco Swing bindings and behavior installations return disposable handles or
are themselves disposable. A \api{BindingScope}, \api{BehaviorScope}, or plain
\api{SubscriptionScope} lets a screen assemble relationships under one owner
and close that owner with the view.
\end{coruscobox}

The strongest habit is local registration. Do not install a listener in one
method and register its cleanup in another method that a reviewer must search
for. If a helper method installs several relationships, return a scope or
disposable that closes all of them. If a presenter activation creates a scope,
pass it to helpers that need to register handles. The code should make
ownership boring enough that there is no temptation to improvise.

Plain Swing listener APIs often make local registration verbose because add and
remove are two separate calls. Wrap the pair immediately. If a document
listener is added to a document, create a subscription whose close method
removes that listener from that document. If a property-change listener is
added to a bean, create a subscription that removes it from the same bean with
the same property name. If a timer is started, register a disposable that stops
it. This tiny adapter is usually better than storing lists of listeners in the
presenter. The presenter should read as ``install this relationship and add its
cleanup'', not as ``remember these five objects so a later method can reverse
engineer what happened.''

\begin{lstlisting}[caption={Wrapping a plain Swing listener as a subscription.}]
DocumentListener listener = new DocumentListenerAdapter() {
    @Override
    public void changedUpdate(DocumentEvent event) {
        presenter.revalidateName();
    }
};
nameField.getDocument().addDocumentListener(listener);

scope.onClose(() -> nameField.getDocument().removeDocumentListener(listener));
\end{lstlisting}

The code is not fancy, but it is honest. The listener and its removal are
adjacent. If a project repeats this pattern often, write a small helper that
returns a \api{Subscription}. The helper should still preserve the ownership
shape: the caller receives a handle and registers it with the scope that owns
the relationship. A helper that hides the handle merely moves the leak to a
different file.

Generated code should follow the same discipline. A generated binding companion
may install several field bindings and validation decorations. It should return
a disposable plan, add everything to a supplied scope, or expose a generated
scope object. It should not install listeners into a view and rely on the view
being garbage-collected. Generation makes repeated wiring cheaper; it does not
change the ownership law. In fact, generated code needs the law more because a
bug in generated cleanup is multiplied across screens.

Scopes also help when cleanup spans different API styles. One child may be a
Corusco subscription. Another may be a \api{Timer} stop callback. Another may
be a table sorter listener removal. Another may be a task handle. The scope
does not care whether each child came from Swing, Corusco core, or application
code. It only knows that they share one lifecycle. This is the practical value
of keeping \api{Disposable} small.

\section{Detachment and reusable state}

Not every lifecycle event is permanent. A tab may deactivate and later
reactivate. A master-detail screen may release cached details when the master
selection changes but keep the model object. A loadable value may drop its
cached result under memory pressure and load again later. A dialog normally
closes permanently, but a workspace panel may detach and return.

\termdef{Detachable}{in Corusco} is a resource that can release cached state
without permanently closing its owner. \api{Detachable.detach} says ``drop the
temporary material you hold now, but remain usable for future activation.''
\api{Disposable.close} says ``this owner is finished; release resources
permanently.'' The distinction prevents two opposite mistakes: keeping heavy
data forever because the model object is reusable, and permanently closing a
model merely because the view is temporarily inactive.

\api{DetachableScope} owns detachable children. Calling \api{detach} releases
cached state in registered children while keeping registrations. Calling
\api{close} detaches once more, clears registrations, and marks the scope
closed. Children detach in reverse registration order. If one child fails, the
scope still attempts to detach the rest and then throws a \api{DetachmentException}
with suppressed failures. If a detachable child is added after close, the child
is detached immediately and not retained.

\begin{lstlisting}[caption={Reusable cache lifecycle.}]
DetachableScope cacheScope = new DetachableScope();

LoadableValue<CustomerDetails> details =
        cacheScope.add(LoadableValue.of(repository::loadDetails));
LoadableList<OrderRow> orders =
        cacheScope.add(LoadableList.of(repository::loadOrders));

cacheScope.detach(); // Release cached details and rows.
cacheScope.close();  // Permanent owner shutdown.
\end{lstlisting}

Detachment is useful when the owner remains meaningful after the current data
is no longer useful. A customer workspace may keep the selected customer id and
command state while detaching the loaded order rows for a hidden detail panel.
A search presenter may keep the query text while detaching cached results when
the screen is deactivated. A map editor may keep layer definitions while
detaching expensive tile caches. In each case the object is not dead; it is
merely no longer actively presenting the cached material.

Detachment is not listener cleanup. If an object subscribed to a long-lived
source, detaching cached data does not remove the subscription unless the
object's own detach contract says so. A detachable loadable value may keep its
ability to load later. A disposable subscription should be closed when the
owner ends. Do not use \api{detach} as a softer word for \api{close}. The two
verbs should have distinct meanings in the project vocabulary.

Activation boundaries become clearer with both verbs available. Opening a
workspace creates a presenter and its permanent scopes. Activating a tab loads
or reattaches detail data. Deactivating the tab detaches heavy cached results
and cancels irrelevant loads. Closing the workspace closes subscriptions,
commands, task services, and detachable scopes. The vocabulary lets the code
say whether a boundary is temporary or final.

A reusable presenter often has two owned scopes: a \api{SubscriptionScope} for
permanent listener relationships and a \api{DetachableScope} for cached data.
The permanent scope closes when the presenter dies. The detachable scope
detaches when the presenter deactivates and closes when the presenter dies.
This pair is more explicit than a single boolean named \api{active}. The
boolean may still exist, but the resource policy has real objects behind it.

\begin{lstlisting}[caption={A reusable presenter has close and detach paths.}]
final class CustomerDetailPresenter implements Disposable, Detachable {
    private final SubscriptionScope subscriptions = new SubscriptionScope();
    private final DetachableScope caches = new DetachableScope();

    @Override
    public void detach() {
        caches.detach();
    }

    @Override
    public void close() {
        try {
            caches.close();
        } finally {
            subscriptions.close();
        }
    }
}
\end{lstlisting}

The order in the example is policy. It detaches caches before closing
subscriptions so cache observers can see a final detached state if needed.
Another application might close subscriptions first. The important fact is that
the order is visible and testable. Hidden cleanup in finalizers or weak
references cannot express such policy.

Detachable state also interacts with background work. A loadable value that is
detached while a load is running should either cancel the load or ignore its
result through generation checks. Otherwise detachment briefly releases memory
and then a late callback repopulates the cache. This is the same stale-result
problem seen in dialogs and search fields; detachment is one more boundary that
invalidates old work.

A master-detail screen gives a good concrete example. The master table
selection is durable screen state. The details for the selected customer are
cached presentation state. When the selection changes, the old details are no
longer useful. The presenter can detach the old detail cache, invalidate the
old detail generation, and start a new load. When the tab deactivates, it may
detach the current detail cache but keep the selected customer id so the tab
can restore context later. When the workspace closes, it closes the detail
presenter permanently. The three events are different, and using both detach
and close lets the code express the difference.

A paint or GIS editor has another pattern. A layer model may remain open for
the lifetime of a document, but rendered tile images, spatial indexes,
selection outlines, and measurement caches may be detachable. When the map
view is hidden, those caches can be dropped. When it is shown again, they can
be rebuilt lazily. The document model should not be closed merely because the
view is hidden. Conversely, the caches should not remain forever merely because
the document remains open. Detachment gives the middle state a name.

Detachment should be observable only when the user or program needs it. A cache
dropping its contents may publish a loadable empty state, a not-loaded state,
or no visible state at all. Choose deliberately. If detaching a hidden tab
causes visible flicker when the tab is shown again, the presenter may need to
distinguish memory-pressure detachment from navigation detachment. The core
concept remains the same, but the presentation policy belongs to the screen.

Tests for detachment should be as concrete as tests for close. Add two tracked
detachable children, call \api{detach}, assert reverse detach order and that
the scope remains open. Call \api{detach} again and decide whether child
detachment should be idempotent. Call \api{close}, assert that children detach
and registrations clear. Add a child after close and assert that it detaches
immediately. These tests prevent the word detach from becoming a vague
ceremony.

\begin{migrationbox}{From caches to detachable state}
When reviewing a presenter with caches, ask which cached values are useful
after deactivation. Keep those as model state. For the rest, introduce
\api{Detachable} objects and collect them in \api{DetachableScope}. Then make
deactivation call \api{detach} and final disposal call \api{close}.
\end{migrationbox}

\section{Task services and callback ownership}

\termdef{UI task}{in Corusco} is work requested by a user-interface owner that
should run away from the caller thread and report a terminal outcome later. The
task may load data, validate a field, save a draft, render a preview, import a
file, or compute a filter. It is called a UI task not because it runs on the
UI thread, but because the result belongs to a UI lifecycle.

\termdef{Task service}{in Corusco} is the lifecycle-owned object that accepts
\api{UiTask} bodies, gives each body a cooperative \api{CancellationToken},
exposes aggregate busy state, and returns a \api{TaskHandle}. The service
separates worker execution from callback delivery. The core service is
toolkit-neutral: its callback executor is supplied by the caller. Swing code
uses a callback executor that returns terminal callbacks and final busy-state
updates to the EDT.

\begin{lstlisting}[caption={Submitting a task through a service.}]
TaskService tasks = TaskService.virtualThreads(EventQueue::invokeLater);

TaskHandle<List<CustomerRow>> handle = tasks.submit(
        cancellation -> {
            cancellation.throwIfCancellationRequested();
            return repository.loadCustomers(cancellation);
        },
        TaskCallbacks.onSuccess(rows::replaceAll));
\end{lstlisting}

The task body runs on a worker. The success callback in this example is posted
through \api{EventQueue.invokeLater}. The task body still must check
cancellation. The callback still must respect the owner lifecycle. A task
service solves the repeated wiring; it does not make asynchronous work
magical.

\api{TaskService.virtualThreads} creates a service backed by Java virtual
threads. A virtual-thread-per-task executor is a good fit for many UI-initiated
blocking operations because each request can have a simple sequential body
without consuming a scarce platform thread while blocked. This does not change
Swing's threading rule. Swing components and Swing-owned models must still be
updated on the EDT. Virtual threads make the worker side easier; they do not
move the UI boundary.

The implementation separates worker ownership from callback ownership. A
service created by the virtual-thread factory owns its worker executor and
shuts it down on close. A service created with explicit executors can cancel
outstanding handles on close while leaving the caller-owned executor service
open. This distinction matters in applications that share a worker pool across
several presenters. Closing one presenter should not shut down the application's
global worker pool unless that presenter owns it.

\termdef{Task handle}{in Corusco} is the object that observes and controls one
submitted task. It exposes task-level busy state, the task's cancellation token,
a completion future, cancellation, and done state. Closing the handle requests
cancellation. It does not close the owning task service. This split lets a
screen cancel one request without shutting down the service that may run later
requests.

\begin{lstlisting}[caption={A screen owns the handles it starts.}]
SubscriptionScope scope = new SubscriptionScope();

TaskHandle<CustomerDetails> handle = taskService.submit(loadDetailsTask,
        TaskCallbacks.onSuccess(details::setValue));

scope.add(handle);
\end{lstlisting}

Registering the handle in a scope states that this particular work belongs to
the screen activation. Closing the scope cancels the handle. If the service
itself also belongs to the activation, register the service too or close it in
the activation cleanup. If the service belongs to a larger presenter, keep the
service open and close only the handles owned by the activation. The distinction
prevents both leaks and overzealous shutdown.

\api{TaskCallbacks} groups terminal callbacks: success, failure, and
cancellation. Only one terminal path should run for a task. A successful body
calls \api{succeeded}; an unexpected failure calls \api{failed}; cooperative
cancellation calls \api{cancelled} and suppresses success and failure. Callback
methods should return quickly because in Swing they normally run on the EDT.
Do not perform long calculations inside terminal callbacks; use them to publish
the already-computed result into models or UI state.

Busy state exists at two levels. Service-level busy means at least one task
owned by the service is running. Handle-level busy means one particular task is
running. A screen-level overlay may observe the service's busy value. A row
progress marker may observe a handle. A Save command may be disabled while the
save handle is busy but leave a Refresh command available if the refresh
service is idle. Use the level that matches the visual policy. Too broad a busy
value freezes unrelated UI; too narrow a value permits conflicting commands.

\begin{figure}[htbp]
\centering
\includegraphics[width=0.88\linewidth]{\modernTaskBoundaryFigure}
\caption{A modern task boundary still returns to Swing ownership: a worker may be virtual, but generation, cancellation, and callback delivery decide whether a result belongs to the current screen.}
\label{fig:lifecycle-modern-task-boundary}
\end{figure}

Figure~\ref{fig:lifecycle-modern-task-boundary} is the whole story in one
picture. Work leaves the EDT because the UI must stay responsive. It carries a
cancellation token because the owner may go away or a newer request may replace
it. It carries a generation token when stale results are possible. Completion
returns through a callback executor because Swing state must be updated on the
EDT. The result is then accepted or ignored according to the lifecycle.

The completion future on a handle represents the final task outcome after the
implementation has performed its documented callback and busy-state delivery.
This is useful in tests and in non-UI coordination. In UI presenters, prefer
callbacks for ordinary model publication because they already arrive on the
configured callback executor. Blocking the EDT waiting for a completion future
is the old SwingWorker mistake in a new costume.

Closing a task service should reject new submissions. That failure is helpful.
Submitting to a closed service means some owner is trying to start work after
its lifetime has ended. Silent acceptance would create a callback with no valid
home. Treat the exception as a wiring bug unless the code is deliberately
racing with shutdown and has an explicit fallback.

A submitted task has a small lifecycle of its own. It is created, marked busy,
scheduled on the worker executor, run with a cancellation token, finished with
success, failure, or cancellation, delivered through the callback executor,
removed from the service's active handle list, and marked not busy. The user
may see only a spinner and later a table update, but the code should preserve
this sequence. Publishing a result before busy state is cleared may be correct
if the UI wants to update data while still showing finalization; publishing it
after busy state is cleared may be correct if the UI wants the result and idle
state together. The service documents its delivery order so presenters and
tests can rely on it.

Task services should not become hidden application globals. A global service
is sometimes useful as a worker substrate, but task handles and callbacks still
belong to UI lifecycles. If a global service stores callbacks directly and the
callbacks capture screens, the global service becomes a leak source. A better
shape is a shared worker executor with presenter-owned task services or
presenter-owned handles. The owner that requested work remains responsible for
its callbacks and cancellation policy.

Failure callbacks deserve the same modeling as success callbacks. A task that
fails during a table refresh may leave existing rows in place and publish a
status error. A task that fails during a first load may publish an empty state
with a retry command. A task that fails during dialog Apply may keep the dialog
open and publish a form-wide problem or workflow error. These policies are not
task-service policies. The service reports the terminal outcome; the owner
decides what failure means in that UI context.

Do not use background tasks to avoid thinking about responsiveness inside the
EDT. If a task callback receives ten thousand rows and then sorts, filters, and
rebuilds a table model on the EDT, the UI may still freeze. Do heavy work in
the task body where possible. Publish compact, ready-to-apply results. Then
perform precise model updates on the EDT. Large data chapters discuss update
precision; the lifecycle point is that a task boundary should reduce EDT work,
not merely postpone it.

\section{Cancellation and stale results}

\termdef{Cancellation}{in Corusco tasks} is a cooperative request for work to
stop. The worker receives a \api{CancellationToken}. It can ask whether
cancellation has been requested and can call \api{throwIfCancellationRequested}
at useful boundaries. The task service also attempts executor-level
cancellation where supported. Neither mechanism proves that every blocking
driver, remote service, or computation stopped instantly.

\begin{pitfallbox}{cancellation is cooperative}
Calling \api{cancel} requests cancellation. It does not prove that a blocking
driver, remote service, or task body has stopped instantly. Task code should
check its token at useful boundaries.
\end{pitfallbox}

\begin{lstlisting}[caption={Cooperative cancellation points.}]
UiTask<List<CustomerRow>> task = cancellation -> {
    cancellation.throwIfCancellationRequested();
    List<CustomerRow> rows = repository.loadCustomers();
    cancellation.throwIfCancellationRequested();
    List<CustomerRow> enriched = enrich(rows, cancellation);
    cancellation.throwIfCancellationRequested();
    return enriched;
};
\end{lstlisting}

Useful boundaries are before expensive work, after blocking calls, inside long
loops, before allocating large result objects, and before publishing a result.
A repository method that accepts the token can check it during paging or
streaming. A method that cannot be cancelled should still be followed by a
token check before its result is used. The user does not care whether the old
request technically finished; the user cares that it no longer changes the
screen after cancellation.

\termdef{Stale result}{in asynchronous UI code} is a result that may be
correct for an older request but no longer belongs to the current presentation
state. A stale result is not necessarily cancelled. A search for \api{Ada} may
finish successfully after the user has already searched for \api{Grace}. A
field validation for an old value may return after a validation for the new
value. A detail load for a previously selected row may finish after the
selection changed. The old result is not wrong in isolation; it is obsolete.

\begin{principlebox}{two gates}
Use cancellation to stop doing useless work. Use generation checks to prevent
obsolete work from becoming visible state.
\end{principlebox}

\api{GenerationCounter} supplies small immutable tokens for stale-result
suppression. Call \api{advance} when scheduling work for the latest input.
Capture the returned generation with the task or callback. Before applying the
result, ask whether the generation is still current or call \api{tryAccept}.
Calling \api{invalidate} advances the counter without scheduling new work,
making outstanding generations stale during close, clear, or reset paths.

\begin{lstlisting}[caption={Suppressing stale search results.}]
GenerationCounter generations = new GenerationCounter();

void search(String text) {
    GenerationCounter.Generation generation = generations.advance();
    taskService.submit(
            cancellation -> repository.search(text, cancellation),
            TaskCallbacks.onSuccess(rows ->
                    generations.tryAccept(generation, rows, tableRows::replaceAll)));
}

void close() {
    generations.invalidate();
}
\end{lstlisting}

The counter is thread-safe because the generation value is atomic. It does not
make the result consumer or the target model thread-safe. If \api{tableRows}
belongs to Swing presentation state, the callback that calls \api{tryAccept}
must already be running on the EDT or must hand accepted work to the EDT. A
generation token decides whether the result is current; it does not decide
which thread may mutate the model.

Cancellation and generations are complementary. Cancelling an old search is
good because it saves work. Keeping a generation check is still necessary
because the old search may finish despite cancellation. Invalidating
generations on close is good because late callbacks cannot publish. Cancelling
handles on close is still necessary because obsolete work should stop when
possible. Treating one mechanism as a replacement for the other is a common
source of stale UI bugs.

\api{AsyncFieldValidation} is a concrete example of both ideas. It subscribes
to a field value, submits validation work to a task service, exposes observable
problem and busy values, and captures a generation for each validation request.
Completions for older values are ignored. Closing the controller invalidates
generations, closes the value subscription, cancels outstanding tasks, clears
tracked handles, and publishes not-busy state. The class is core-level and
Swing-free; Swing callers choose a task service whose callbacks return to the
EDT.

The async validator's failure policy is also important. Unexpected validator
failures are converted into a field-targeted error problem with a stable
problem code. This keeps the form in a coherent state: the user sees that the
asynchronous validation failed, OK can become non-committable if errors block
commit, and tests can assert the problem code. A thrown exception does not
escape into a worker log while the dialog pretends the field is valid.

When a field value changes rapidly, busy state should represent the current
generation. Older generations should not clear busy for newer work. Otherwise
the user sees a flicker: first request starts, second request starts, first
request finishes and clears busy, second request still runs silently. The
current-generation check in async validation prevents this. The same principle
applies to table loads, preview renders, and search requests.

Generation counters should not be persisted or compared by raw number. The
numeric value is meaningful only to the counter that produced it. Treat the
token as an opaque ownership proof: it says ``this result belongs to the
request that was current when the token was captured, unless the counter has
advanced.'' This keeps code from growing accidental ordering logic around
generation values.

There are several common stale-result shapes. Search-as-you-type creates many
requests for a changing text value. Master-detail loading creates requests for
changing selection. Asynchronous validation creates requests for changing field
values. Preview rendering creates requests for changing options. Import
analysis creates requests for changing files. In every case the current input
can change while old work is still running. A generation counter is not a
special validation tool; it is a general ownership token for request streams.

Sometimes the stale-result boundary is not user input but lifecycle. Closing a
dialog invalidates all outstanding dialog requests. Deactivating a tab may
invalidate detail loads. Clearing a form may invalidate validation work.
Switching projects may invalidate all workspace loads. The method name
\api{invalidate} is useful because it says that no old result should be
accepted even if no newer request is started immediately. This is different
from \api{advance} as a scheduling phrase, although the implementation advances
the counter in both cases.

Stale suppression should be as close as possible to publication. A worker can
check cancellation before returning, but it usually cannot know whether its
result is still current by the time the callback runs. The callback is the last
gate before state changes. Put the generation check there. If the callback
delegates to another method, make the generation part of that method's
contract or keep the check before the call. A generation captured at scheduling
time and forgotten by callback time has no value.

When stale results are ignored, consider whether diagnostics should record the
fact. In ordinary typing, ignored stale validations are normal and should not
be logged as warnings. During shutdown, ignored callbacks are also normal. But
if a long-running refresh is routinely made stale by another refresh, the UI
may be scheduling too aggressively. Counters can support quiet correctness
while metrics or debug logs support performance tuning.

\section{The Swing boundary and busy UI}

Swing's Event Dispatch Thread is the boundary that every UI task must respect.
Task bodies run away from the EDT. Terminal callbacks that update Swing
components, Swing models, or Swing-owned Corusco values must run on the EDT.
\api{SwingTaskServices.virtualThreads} supplies the common Swing choice: a
virtual-thread worker service whose terminal callbacks and final busy-state
updates are delivered through \api{SwingEdt.executor}. Presenters can submit
work without writing \api{SwingUtilities.invokeLater} at every call site.

\begin{lstlisting}[caption={Swing task service for a presenter or dialog.}]
TaskService tasks = SwingTaskServices.virtualThreads();
FormDialogLifecycle lifecycle = FormDialogLifecycle.of(dialog);
lifecycle.addTaskService(tasks);
\end{lstlisting}

The returned service is a lifecycle object. It should be closed by the owning
application, presenter, screen activation, or dialog lifecycle. A service
owned by a dialog should end with the dialog. A service owned by a workspace
presenter may survive view deactivation if that is the presenter's policy. A
global service should be rare and should not retain screen callbacks without
handles or generation checks.

Busy overlays are a presentation of busy state, not a substitute for task
ownership. A \api{BusyOverlayBinding} can observe a readable busy value and
update a \api{JLayer} with \api{BusyOverlayLayerUI}. The overlay may paint a
translucent cover and consume input while busy. The binding itself touches
Swing components and must be closed with the view. If the overlay remains
installed after the view closes, the visual state was not the only leak; the
binding still owns listeners and component references.

Command enablement should follow the same modeled state. A Refresh command may
be disabled while a refresh handle is busy. A Save command may be disabled
while validation is pending or while a save is running. A Cancel current task
command may be enabled only while a handle is busy and not done. These are
ordinary derived values. They should be bound to commands and components, not
hand-updated in every task callback.

Dialogs make the Swing boundary especially visible. A dialog may start a
remote validation task. The user may press Cancel. The lifecycle closes the
validation binding, cancels the handle, invalidates the generation, and closes
the controller. A late callback returns on the EDT; it sees that its generation
is stale or its owner is closed and does not publish. Without this chain, the
closed dialog may still mutate a form model, update a summary label, or
re-enable OK after it no longer exists.

Full screens have a wider boundary. A customer workspace may keep a task
service for the whole workspace, but each selected customer detail load has a
request generation. Changing selection cancels the previous detail handle when
possible, advances the generation, shows detail busy state, and accepts only
the current result. Closing the workspace closes the service or all owned
handles according to ownership. The same pattern scales from dialogs to
workspaces because the concepts are not visual.

Status text should be modeled, not scattered. A task may publish ``Loading
customers'', ``Loaded 200 rows'', ``Cancelled'', or ``Failed: network timeout''
through callbacks. The status label observes the modeled status. A failure
callback may also publish a problem, enable retry, and record diagnostics. If
status text is assigned directly from worker code, threading and stale-result
policy become invisible.

Progress deserves the same care. An indeterminate progress bar can observe a
busy value. Determinate progress needs a modeled progress value. If progress
belongs to one handle, close the handle and its progress subscription together.
If progress belongs to a service with several concurrent tasks, define how
aggregate progress is computed or avoid pretending it is determinate. False
precision is worse than an honest busy indicator.

Do not block the EDT waiting for background work. A modal dialog is already a
structured waiting experience for the user; blocking the EDT inside that modal
experience freezes painting, input, cancellation, and progress. If OK needs a
server save before closing, start a task, put the dialog in busy state, keep
Cancel or Stop policy explicit, and close only after the EDT callback reports
success. If OK can return a result and the surrounding screen can save, close
the dialog and let the surrounding workflow own the save task. Both designs are
valid; blocking the EDT is not.

The Swing boundary also affects model classes that are not Swing components
but are owned by Swing presentation. An observable list feeding a table model,
a form model bound to fields, or a command model adapted to buttons may be
toolkit-neutral in package name but still confined to the EDT by ownership. If
the model is read and written only by the UI, treat it as EDT-owned. A worker
should return immutable data or a batch description; the EDT callback applies
it. If a model is truly thread-safe, document that fact separately. Do not infer
thread safety from the absence of Swing imports.

Busy overlays should have input policy. An overlay that paints busy state but
lets all input through may be suitable for a nonblocking refresh. An overlay
that consumes input may be suitable when edits would conflict with a running
operation. A Cancel button may need to remain above or outside the overlay.
These are design choices connected to task ownership. If the task can be
cancelled, the UI should expose a cancellation route. If it cannot be
cancelled, the UI should avoid pretending that pressing Cancel will stop the
work instantly.

Command state can make busy policy precise. Instead of disabling an entire
panel, disable Save while saving, disable Refresh while refreshing, keep Help
enabled, and keep Cancel current task enabled while the relevant handle is
busy. A composed command enablement value can depend on validation problems,
dirty state, busy state, and lifecycle state. Corusco's value and command
models make these dependencies explicit; Swing buttons merely render the
command state.

For screenshot generation, busy state should be deterministic. The book's
figures use seeded data and fixed dimensions, but the same principle applies
to application screenshot tests. Drive the model into a busy state, render the
view, then drive it into a loaded or failed state. Do not sleep and hope a real
background task happens to be at the right point. Deterministic screenshots are
also better documentation because they show named states rather than timing
accidents.

\section{Testing lifecycle and tasks}

Lifecycle behavior deserves tests. Not every screen needs a heap test, but
every screen pattern should have cleanup tests. Close a scope and assert that
later value changes do not update the view. Close a task handle and assert that
the token is cancelled. Invalidate a generation and assert that the old result
is ignored. Close a lifecycle and assert reverse cleanup order. These tests are
small because the contracts are small.

\begin{lstlisting}[caption={Testing a closed subscription scope.}]
SimpleValue<String> value = SimpleValue.of("first");
SubscriptionScope scope = new SubscriptionScope();
List<String> seen = new ArrayList<>();

scope.add(value.subscribe(event -> seen.add(event.newValue())));
scope.close();
value.setValue("second");

assertThat(seen).isEmpty();
\end{lstlisting}

The test above is more valuable than its size suggests. It proves that a
closed owner does not receive later model updates. A similar test can assert
that late registration closes immediately: close a scope, add a tracked
disposable, and verify that its close method ran. Another test can register
three disposables and assert reverse close order. If the project has a custom
binding scope, give it the same tests.

Task tests should control time. Use a direct callback executor for core tests
when no Swing boundary is involved. Use a latch or manually controlled task
body to hold a task in the busy state. Assert service busy, handle busy,
callback delivery, completion future behavior, cancellation, and close
semantics. For Swing tests, use \api{SwingEdt.runAndWait} or a Swing-aware test
helper so callback assertions happen on the EDT.

\begin{lstlisting}[caption={A deterministic core task scenario.}]
CountDownLatch release = new CountDownLatch(1);
try (TaskService service = TaskService.virtualThreads(Runnable::run)) {
    TaskHandle<String> handle = service.submit(
            cancellation -> {
                cancellation.throwIfCancellationRequested();
                release.await();
                return "loaded";
            },
            TaskCallbacks.onSuccess(values::add));

    assertThat(handle.busy().value()).isTrue();
    release.countDown();
    assertThat(handle.completion().join()).isEqualTo("loaded");
}
\end{lstlisting}

Stale-result tests do not need real threads. Capture two generations. Try to
accept the older result after advancing the counter. Assert that only the newer
result is accepted. Then invalidate the counter and assert that a previously
current token becomes stale. This proves the callback gate independently from
worker scheduling. A separate integration test can prove that the task service
delivers callbacks through the expected executor.

Cancellation tests should distinguish request from completion. After calling
\api{cancel}, the token should report cancellation. The completion may be
cancelled after callback delivery according to the service contract. If the
task body checks the token, assert that it stops at the expected boundary. If
the body blocks in a non-cancellable call, assert that a late result is ignored
by generation or lifecycle checks. Do not write tests that assume every
blocking operation stops instantly.

Cleanup failure tests are worth having for scopes used across the project.
Register one disposable that throws and another that records closure. Close the
scope and assert that the recording disposable still closed and the thrown
exception appears as suppressed inside the aggregate exception. This prevents a
future refactoring from accidentally changing cleanup from best-effort to
stop-on-first-failure.

Screenshot tests can reveal lifecycle policy when they are model-driven. A
busy overlay screenshot should come from a busy value. A disabled command
should come from command state. A failure banner should come from task failure
state. If the screenshot can be produced only by manually painting a component
without the underlying model state, the example is less useful. Visual evidence
and behavioral tests should describe the same relationships.

Fake services are useful when they preserve the public contract. A fake task
service for presenter tests can record submitted tasks, expose a busy value,
and let the test complete, fail, or cancel a handle manually. It should still
return a \api{TaskHandle}; it should still support cancellation; it should
still deliver callbacks through the chosen test executor. A fake that simply
calls the success callback immediately is convenient, but it cannot test busy
state, cancellation, or stale results. Use the simplest fake that still covers
the lifecycle rule under test.

EDT tests should avoid nested modal waits. Test controllers, presenters, and
bindings directly on the EDT. For a dialog, drive the \api{FormDialog}
controller and lifecycle without showing a native modal peer. For a busy
overlay, render the layer with a controlled busy value. For task callbacks,
use \api{SwingEdt.runAndWait} around state assertions or a test utility that
flushes posted EDT work. Native modal integration can have a small smoke test,
but most lifecycle behavior should be tested without clicking windows.

Heap tests are a last line of defense, not the primary design tool. They can
prove that a screen becomes collectible after close under controlled
conditions, but they are often slow and sensitive to test infrastructure.
Prefer direct ownership tests first: listener count goes to zero, scope closes
children, callbacks do not run after close, task handles cancel, generations
reject stale results. When these tests pass, heap tests become confirmation
rather than discovery.

Logging should identify ownership. If a task fails after its owner has closed,
that may be a debug-level stale result, not an application error. If a task is
submitted to a closed service, that is probably a wiring error. If cleanup
throws, the aggregate exception should retain suppressed failures so the
developer sees all broken cleanup handles. Good diagnostics use lifecycle
vocabulary; poor diagnostics merely say that a null component was updated.

\begin{examplebox}{lifecycle}
\lifecycleExample{} demonstrates deterministic cleanup of listener-style
registrations. \asyncTaskExample{} demonstrates core task execution and busy
state. \generationCounterExample{} demonstrates stale-result suppression
without real threads, so the ownership rule is easy to see.
\end{examplebox}

\section{Worked review of a customer workspace}

Consider a customer workspace with a search field, a table of results, a detail
panel for the selected customer, an Edit dialog, and a status bar. The screen
is ordinary enough to be useful and rich enough to expose lifecycle questions.
The search field starts background requests. The table observes an observable
list. The detail panel loads data for the selected row. The Edit dialog has its
own modal lifecycle. The status bar observes task and command state. The
workspace may be closed and reopened during one application session.

The first review pass marks relationships installed during workspace
activation. The search document listener, table selection listener, list
model adapter, command adapters, status binding, busy overlay binding, and
task service all need owners. Some belong to the whole workspace activation.
Some belong to the current detail selection. Some belong only to the Edit
dialog. Marking them by owner is more useful than listing them by class.

The workspace presenter creates a \api{SubscriptionScope} for activation. It
registers the search listener, table binding, status binding, command
adapters, and busy overlay binding. It also creates a \api{TaskService} for
workspace tasks and registers it if the service is not shared elsewhere. The
activation close method closes the scope. The view disposal method calls the
presenter close method. This is not complicated code, but it must be present.

Search requests use a \api{GenerationCounter}. Each change that starts a
search advances the generation and captures the token. The old search handle
may be cancelled if still running. The task callback returns on the EDT. It
tries to accept the result with the captured generation before replacing table
rows. Clearing the search field invalidates generations even if no new search
is scheduled. Closing the workspace invalidates generations and closes handles
or the service. Thus old results cannot repopulate the table after close or
after a newer search.

The detail panel has a different boundary: selection lifetime. Selecting a row
starts a detail load for that row. Changing selection cancels or invalidates
the old detail load and detaches cached detail state if the application wants
to release memory. The workspace may keep the last selected id as durable
state, while detaching the heavy order list and preview image. A
\api{DetachableScope} fits this boundary better than closing the whole
workspace presenter.

The Edit dialog creates a new modal lifecycle. It may use the same repository
service, but its form bindings, validation binding, keyboard binding, dialog
task service or handles, dirty confirmation, and controller belong to the
dialog activation. Closing the workspace while the dialog is open should close
the dialog lifecycle. Cancelling the dialog should not close the whole
workspace task service unless the dialog owns that service. These ownership
lines prevent modal cleanup from damaging the surrounding screen.

Now follow one user gesture. The user types \api{Ada} in the search field. The
document listener, owned by the workspace activation scope, asks the presenter
to search. The presenter advances the search generation, cancels the previous
search handle if one exists, submits a new task through the workspace task
service, and stores the handle in the activation scope or in a small
replaceable handle owner. The service busy value turns true, the overlay
binding shows the busy layer, and the Refresh command disables. The user types
\api{Ada L}. The presenter repeats the sequence. When the first task returns,
its callback runs on the EDT but fails the generation check and publishes
nothing. When the second task returns, its generation is current, so the table
rows update and busy state clears.

Follow another gesture. The user selects a customer row. The selection listener
belongs to the workspace activation, but the detail load belongs to the
selection. The presenter invalidates the previous detail generation, detaches
old heavy detail caches, and submits a new detail load. The detail panel shows
busy state without blocking the search table. If the user selects a different
row before the load finishes, the old result is ignored. If the user closes the
workspace, the activation scope closes the selection listener, cancels the
detail handle, closes or detaches detail caches according to policy, and
invalidates outstanding generations. The late detail result has no valid path
back to the closed view.

Follow the Edit dialog. The user opens it from the selected row. The dialog
creates a fresh form model, root panel, controller, lifecycle, validation
binding, keyboard binding, and perhaps a dialog-owned task service for
asynchronous validation. The dialog lifecycle owns these objects. The
workspace may remain open behind it, but the dialog does not borrow the
workspace activation scope for modal resources. If the user presses Cancel,
the dialog cancels or invalidates dialog work and closes dialog resources. The
workspace still owns its search and detail tasks. If the workspace closes, it
also closes any open dialog lifecycle because the modal child can no longer
outlive its owner.

This walkthrough is verbose on purpose. Real bugs often occur at the boundary
between two lifecycles: search vs detail, detail vs dialog, dialog vs
workspace, workspace vs application shutdown. A diagram or ledger that makes
the boundaries visible prevents accidental ownership transfer. Without that,
developers tend to choose whatever scope is in reach, and a dialog task ends up
owned by a workspace or a workspace listener ends up owned by a dialog.

Failure paths need explicit owners. If a search fails, the workspace status
model receives the failure and the table remains in its previous or empty
state according to policy. If a detail load fails, the detail panel displays a
detail-level failure and the search table remains usable. If an Edit dialog
save fails while the dialog stays open, the dialog status or form problem state
receives the failure. The same exception class may be involved in all three
paths; the owner decides the presentation.

Shutdown has different rules from ordinary navigation. During normal workspace
close, user policy may ask about dirty edits and may cancel background work.
During application shutdown, windows may already be disappearing and user
prompts may be limited. The application should define which task services close
first, whether outstanding saves are awaited or cancelled, and which callbacks
are ignored after shutdown begins. Corusco's scopes and task services give
places to implement that policy; they do not invent the policy for the
application.

The review table for this workspace is concrete:

\begin{center}
\textbf{Lifecycle review table.}
\end{center}

\noindent\begin{tabularx}{\linewidth}{@{}p{0.34\linewidth}X@{}}
\toprule
Observation & Review question\\
\midrule
Validation runs twice after reopening & Was the old value or document subscription closed?\\
Closed dialog updates status & Does the task callback check lifecycle or generation before publishing?\\
Search results revert to older text & Does each request capture and verify a generation token?\\
Cancel does not stop loading & Are task handles registered and cancelled at the right boundary?\\
Busy overlay remains after close & Is the busy binding closed with the view?\\
Detail rows reload after tab deactivation & Should the detail cache detach or should the load be cancelled?\\
Closing one panel stops all background work & Does the panel own the task service, or only its handles?\\
Cleanup exception leaves listeners alive & Does the scope continue closing remaining children?\\
\bottomrule
\end{tabularx}

The questions are not theoretical. They are the bugs that appear in long-lived
desktop tools. A one-shot demo can ignore them and exit. A production Swing
application stays open all day. Time magnifies every unclear ownership rule.

\section{Operational notes}

Lifecycle policy should be visible in diagnostics. When a task is submitted,
debug logging can include the owner, operation name, and generation token
identity without exposing domain data. When a task result is ignored because it
is stale, that is usually a trace-level fact, not an error. When code submits a
task to a closed service, that is usually a warning or exception because a dead
owner is trying to start new work. When cleanup fails, the aggregate exception
should retain suppressed failures so the developer sees every broken cleanup
handle. Diagnostics that use lifecycle vocabulary are easier to act on than a
late \api{NullPointerException} in a callback.

Memory investigation also becomes simpler when ownership is explicit. If a
heap dump shows old customer panels retained, inspect long-lived listener
sources, task services, command registries, and static caches. The ownership
ledger tells which relationships should have been closed. If the retained path
goes through a subscription, find the scope that should have owned it. If the
retained path goes through a task callback, find the handle or service that
should have been cancelled. If the retained path goes through a cache, ask
whether the cache should detach at deactivation. The heap tool tells what is
still reachable; the lifecycle design tells why that reachability is wrong.

Performance investigation follows the same pattern. Duplicate callbacks often
look like slow validation or slow table updates. The real cause may be three
live subscriptions all responding to the same field change. A busy overlay that
stays visible may look like a painting defect; the real cause may be a task
handle whose busy value never cleared because completion was delivered to a
closed executor. A UI freeze may look like a background-task problem; the real
cause may be a success callback doing heavy work on the EDT. Naming the owner,
callback executor, and generation makes these problems easier to separate.

Application shutdown deserves a written sequence. Top-level windows stop
accepting new user work. Dialog lifecycles close or cancel according to policy.
Screen activation scopes close. Presenter-owned task services cancel
outstanding handles. Shared services stop accepting new submissions and then
shut down if the application owns them. Final callbacks after shutdown begin
are either ignored by invalidated generations or routed to shutdown diagnostics
instead of UI state. A shutdown sequence does not need to be elaborate, but it
must be deliberate; otherwise every window invents a different interpretation
of closing.

The distinction between close and dispose should remain clear in operations
documentation. Closing a Corusco scope releases Java relationships. Disposing a
Swing window releases the native peer. Cancelling a task requests worker
cooperation. Invalidating a generation prevents old results from publishing.
Detaching a cache releases reusable cached state. These operations often occur
together, but they are not synonyms. A support note that says ``the dialog is
disposed'' does not prove its validation task was cancelled. A code review that
says ``the task is cancelled'' does not prove old callbacks are unable to
publish. Use the exact verb.

Instrumentation should not become another leak. A diagnostic listener attached
to a task service needs an owner. A metrics observer attached to an observable
list needs a subscription. A debug overlay installed on a \api{JLayer} needs a
binding handle. Development-only tools are especially prone to unclear
ownership because they are added after the main design is finished. Treat them
like production relationships. If they observe a screen, close them with the
screen. If they observe the application, make sure they do not capture screens.

Code generation should emit lifecycle hints where possible. A generated form
companion can name the binding scope it expects. A generated table companion
can expose a single install method that returns a disposable plan. A generated
action companion can document whether command adapters are owned by the view or
by a longer-lived command registry. A generated validation companion can use
generation counters for asynchronous validators and expose the controller as a
disposable. Generated code is best when it reduces repeated ownership wiring
without hiding ownership.

Finally, keep the human debugging story in mind. A developer called to fix a
stale result bug should be able to answer four questions quickly: who started
the work, who can cancel it, how does the callback know it is still current,
and who closes the callback relationship? If the answer requires reading
twenty anonymous listeners across a dialog and presenter, the architecture is
too implicit. If the answer is ``the customer workspace activation scope owns
the handle, the search generation gates publication, and callbacks return on
the EDT,'' the system is easier to maintain.

\section{Exercises and checklist}

\begin{exercisebox}
\begin{enumerate}
\item Pick one existing screen and mark every listener, subscription, timer,
task, model adapter, keyboard binding, and overlay it installs. For each mark,
write the owner that closes it.
\item Replace two manually stored listener-removal lambdas with a
\api{SubscriptionScope}. Add a test for reverse close order.
\item Add a late-registration test to a custom scope or presenter helper. Close
the owner first, then register a tracked disposable and assert that it closes
immediately.
\item Identify one cache that should survive object lifetime but not activation
lifetime. Implement \api{Detachable} and collect it in a \api{DetachableScope}.
\item Convert one \api{SwingWorker} use to a \api{TaskService} and
\api{TaskHandle}. Preserve callback delivery on the EDT.
\item Add cooperative cancellation checks before expensive work, after a
blocking call, and before publishing a result.
\item Add a \api{GenerationCounter} to a search, validation, or detail-load
path. Prove that an older result cannot replace a newer one.
\item Bind a service-level busy value to a screen overlay and a handle-level
busy value to a command. Explain why the levels differ.
\item Test that closing a dialog lifecycle cancels outstanding task handles and
prevents late callbacks from updating form state.
\item Review a shutdown path and decide which services close, which handles
cancel, and which callbacks are ignored after shutdown begins.
\end{enumerate}
\end{exercisebox}

\begin{center}
\textbf{Lifecycle checklist.}
\end{center}

\begin{itemize}
\item Every installed callback has a visible owner.
\item Listener and binding cleanup handles are registered near installation.
\item \api{SubscriptionScope} closes screen or dialog relationships.
\item Late registrations fail closed instead of creating fresh leaks.
\item Cleanup failures do not prevent remaining children from closing.
\item \api{detach} is used for reusable cached state, not as a synonym for close.
\item Task services have an owner and reject submissions after close.
\item Task handles are registered with the lifecycle that requested the work.
\item Cancellation tokens are checked at useful worker boundaries.
\item Generation tokens suppress obsolete callback results.
\item Swing callbacks that mutate UI state return to the EDT.
\item Busy state is modeled at the correct service or handle level.
\item Dialogs, tabs, selections, and application shutdown have distinct boundaries.
\item Tests cover close, detach, cancellation, stale results, and late callbacks.
\item Example screenshots are backed by modeled busy and task state.
\end{itemize}

Lifecycle code should be boring and visible. Clever cleanup hidden in
finalizers, weak listeners used without policy, callbacks that check random
booleans, and tasks that publish into whatever object they captured are harder
to review than scopes and handles. Use the small vocabulary consistently:
open, activate, detach, cancel, invalidate, close, and dispose. When those
words have stable meanings, a large Swing application becomes much easier to
keep alive for a long time.
